% ── 全文通用形态（Mrite 高教社杯规范）──
% 正文一律自然段；禁分点符号、禁 	extbf（摘要“针对问题X”与问题重述“问题N：”除外）。
% 跨页表用 longtable，短表用 tabularx；列宽比例总和 = 1.04 − 0.04 × N（N 为列数），格内文字不换行。
% 图宽 0.8	extwidth，表宽 	extwidth；图题交给 \caption，绘图代码不写 set_title()；图表必须在正文被引用。
% 章与章、节与节之间一律不加 
ewpage。
% ── 本节合同 ──
% 优点写“做到了什么 + 证据在哪”，不写空泛褒词。
% 缺点按“缺陷 — 影响 — 改进”三段写，且必须与降级声明一致：凡被降级的结论，此处要写明实际强度等级。
%
\section{模型的评价}

\subsection{模型的优点}
\begin{itemize}[itemindent=2em]
\item \textbf{优点1：}
\item \textbf{优点2：}
\item \textbf{优点3：}
\end{itemize}

\subsection{模型的缺点}
\begin{itemize}[itemindent=2em]
\item \textbf{缺点1：}
\item \textbf{缺点2：}
\end{itemize}
